% Level 3 answer key, Week 7 -- Grids, Recursion and Fractal Art
\answersweek{7}{Grids, Recursion and Fractal Art}{Answers for workbook chapter 7 and homework 7.}

\section{Workbook: questions in the text}
\begin{answer}[Predict: \texttt{stars(3)}]
Three lines: \texttt{***}, then \texttt{**}, then \texttt{*}. The base case \py|n == 0| stops it before an empty line is printed. Accept any answer that shows the lines getting shorter and stopping.
\end{answer}
\begin{answer}[Your turn: grid cells]
\texttt{grid[0][3]} is \texttt{\textasciitilde} (water, top-right corner); \texttt{grid[2][1]} is \texttt{.} (land); the mountain \texttt{\textasciicircum} is at \texttt{grid[2][2]}. A student who answers \texttt{grid[2][1]} for the mountain has swapped the order or miscounted from 1: point at the picture and say \enquote{down to row 2 first, then across to column 2, both counted from 0.}
\end{answer}

\section{Exercise 7.1: Recursion warm-up}
\pyfile{../exercises/week07/ex-countdown-sum/solution.py}
\begin{answer}[What to look for]
Base case of \texttt{sum\_to}: \py|if n <= 0: return 0| (\py|n == 0| is also fine). The test \enquote{sum to calls itself} fails if a student used a loop: praise the working loop, then ask them to rewrite it recursively -- that is the point of the exercise. For \texttt{reverse}, another correct recursive case is \py|return reverse(word[1:]) + word[0]|.
\end{answer}

\section{Exercise 7.2: Fractal tree}
\pyfile{../exercises/week07/ex-fractal-tree/solution.py}
\begin{answer}[Table and pattern]
Depth 1, 2, 3, 4, 5 $\rightarrow$ 1, 3, \textbf{7, 15, 31} lines (the depth-5 tree is the one on the right of the screen). Pattern: each depth doubles the previous number and adds 1 (the new trunk); it is always one less than a power of 2 ($2^{\text{depth}} - 1$). Accept \enquote{it doubles (plus one)}. The tests also accept trees that walk back with the pen down, and any \texttt{ANGLE}/\texttt{SHRINK} values.
\end{answer}

\section{Exercise 7.3: Paint bucket}
\pyfile[firstline=53,lastline=92]{../exercises/week07/ex-paint-bucket/solution.py}
\begin{answer}[What to look for]
The paint stops at the \textbf{edge of the grid}, at \textbf{walls} (\texttt{\#}) and at any cell that is \textbf{not the old color} (including cells already painted). Those are the base cases. The four recursive calls can be in any order. A missing \enquote{not the old color} check gives \texttt{RecursionError} -- the planned bug hunt. The drawing and click code (from the \texttt{drawing} comment onwards in \texttt{solution.py}) is given to students unchanged.
\end{answer}

\section{Stretch: BFS with a queue}
\pyfile[firstline=95,lastline=115]{../exercises/week07/ex-paint-bucket/solution.py}
\begin{answer}[What to look for (optional stretch)]
Workbook question: in an empty 4-by-4 grid started at the top-left, the \textbf{bottom-right} square is visited last (it is 6 steps away). The loop must (1) \texttt{popleft()} from a \texttt{deque}, (2) \texttt{append} neighbors, and (3) mark a square \emph{when it joins the queue} so that it is never added twice. Forgetting the mark does not give a \texttt{RecursionError} but a loop that never ends (the queue keeps refilling). The result must equal \texttt{flood\_fill}: same painted squares. Using \texttt{pop()} instead of \texttt{popleft()} still fills correctly, but it visits in depth-first order, not rings.
\end{answer}

\section{Homework 7}
\subsection{Core: Treasure map}
\pyfile[firstline=43,lastline=63]{../exercises/week07/hw-treasure-map/solution.py}
\begin{answer}[Core]
A corner cell has \textbf{2} neighbors (an edge cell 3); a middle cell has \textbf{4}. The program prints \texttt{Start is at (2, 3)} and \texttt{Next to the start: [(1, 3), (3, 3), (2, 2), (2, 4)]}, and the words above the map say \enquote{Treasure on this map: 2}. A loop with \py|if cell == tile| instead of \py|.count()| is equally good.
\end{answer}

\subsection{Challenge: Koch snowflake}
\pyfile[firstline=19,lastline=43]{../exercises/week07/hw-koch-snowflake/solution.py}
\begin{answer}[Challenge]
Depth 2 = \textbf{16} lines per side, depth 3 = \textbf{64} (each level multiplies by 4). Base case: depth 0 draws one straight line. If a student turns \texttt{left(120)} in \texttt{snowflake()}, the bumps point inwards -- an \enquote{anti-snowflake}, also a nice result.
\end{answer}

\begin{answer}[Challenge: BFS rings (stretch)]
Visiting order (up, down, left, right):
\begin{center}
\begin{tabular}{ccc}
1 & 3 & 6\\ 2 & 5 & 8\\ 4 & 7 & 9\\
\end{tabular}
\end{center}
The homework starts in the \textbf{middle} square (the workbook example starts in a corner, so students cannot copy it). Visiting order (up, down, left, right):
\begin{center}
\begin{tabular}{ccc}
6 & 2 & 7\\ 4 & 1 & 5\\ 8 & 3 & 9\\
\end{tabular}
\end{center}
The highest number is \textbf{9}, the \textbf{bottom-right} square, \textbf{2} steps from the start. Distances: the four side squares (2, 3, 4, 5) are 1 step away; the four corners (6, 7, 8, 9) are 2 steps away. Accept a different order inside a ring if the student states a different neighbor order, but every ring must be finished before the next one starts: all of 2--5 must be next to the start and all of 6--9 must be corners.
\end{answer}

\begin{answer}[Creative: my own fractal]
Look for: a base case they can point to, a visible repeating structure, and a \texttt{\# Based on: ...} credit line. Any recursive pattern counts, even a simple one.
\end{answer}
